\documentclass[italian,12pt]{article}

%--------------variabili------------------%
\def\Title{Glossario}
\def\Author{Code7Crusaders}
\def\Version{v0.1}
%-----------------------------------------%

%pacchetti extra da scaricare dblfloatfix, fancyhdr
\usepackage{fancyhdr}%creazione header-footer
\usepackage{graphicx} %serve per inserire immagini
%\usepackage{dblfloatfix} %serve per posizionare gli elementi dove si vuole
\usepackage[hidelinks]{hyperref} %serve per i link
\usepackage{tikz}
\usepackage{tgadventor} % font
\usepackage[useregional=numeric,showseconds=true,showzone=false]{datetime2}
\usepackage{caption}
\usepackage[automake,section=section,toc]{glossaries-extra}
\usepackage{geometry}
\usepackage{setspace}
\usepackage{eurosym}
\usepackage[italian]{babel}
\usepackage[hidelinks]{hyperref}

\linespread{1.2}
\captionsetup[table]{labelformat=empty}
\geometry{a4paper, margin=1in}

\renewcommand{\contentsname}{Indice}%imposto il nome dell'indice
\renewcommand\familydefault{\sfdefault}

\glstoctrue

\makeglossaries

% Definizioni delle voci
%-----------------------TUTORIAL PER AGGIUNGERE VOCI-----------------------------------%
%\newglossaryentry{NOME}{
%    name=NOME, %IMPORTANTE NON MANCARE LA VIRGOLA
%   description={METTERE QUI UNA DESCRIZIONE}
%}
%%-------------Inserire termini da qui
\newglossaryentry{Glossario}{
    name=Glossario,
    description={Elenco strutturato di termini tecnici o specializzati, ognuno corredato dalla propria definizione o spiegazione. Questo strumento aiuta a migliorare la comunicazione tra le varie parti coinvolte in un progetto, riducendo le ambiguità e garantendo una comprensione condivisa dei termini utilizzati in un determinato contesto.}
}

\newglossaryentry{Norme di Progetto}{
    name={Norme di Progetto},
    description={Regole e linee guida stabilite all'interno di un progetto per garantire coerenza e qualità nelle attività svolte. Definiscono standard e procedure, come documentazione, gestione delle versioni e criteri di codifica, per assicurare uniformità nell'approccio e nel risultato finale.}
}

\newglossaryentry{Piano di Progetto}{
    name={Piano di Progetto},
    description={Documento formale che delinea in dettaglio la pianificazione, la esecuzione, il monitoraggio e il controllo di tutte le attività coinvolte nella realizzazione di un progetto. Questo documento fornisce una roadmap chiara e organizzata, comprensiva di obiettivi, risorse, scadenze e strategie di gestione dei rischi.}
}

\newglossaryentry{Piano di Qualifica}{
    name={Piano di Qualifica},
    description={Documento che stabilisce gli standard di qualità, i processi e le attività di testing che saranno implementati durante lo sviluppo di un progetto. Contiene una descrizione dettagliata delle strategie di testing, delle metriche di valutazione e dei criteri di accettazione del prodotto finale.}
}

\newglossarystyle{myaltlistgroup}{%
    \setglossarystyle{altlistgroup}%
    \renewcommand*{\glsgroupheading}[1]{%
        \pagebreak%
        \item\par\bigskip\textbf{\Large \glsgetgrouptitle{##1}}\par\medskip % Titolo della lettera in grande
        \item\par\noindent\hrule\par\medskip % Linea di separazione
    }%
}


\glsaddallunused
%------------------------Inizio del documento--------------
\begin{document}

% Sezione del titolo
\begin{titlepage}

    \AddToHookNext{shipout/background}{
    \begin{tikzpicture}[remember picture,overlay]
    \node at (current page.center) {
    \includegraphics{../../img/background.png}
    };
    \end{tikzpicture}
    }

    \centering
    \vspace*{2cm}
    
    \includegraphics[width=0.3\textwidth]{../../img/logo/7Crusaders_logo.png} % Aggiungi il logo qui
    \vspace{1cm}
    
    {\Huge \textbf{Code7Crusaders}}\\
    \vspace{0.5cm}
    {\Large Software Development Team}\\
    \vspace{2cm}
    
    \large \textbf{Glossario}
    \vspace{3.9cm}

    \textbf{Membri del Team:}\\
    Enrico Cotti Cottini, Gabriele Di Pietro, Tommaso Diviesti \\
    Francesco Lapenna, Matthew Pan, Eddy Pinarello, Filippo Rizzolo \\
    \vspace{0.5cm}
    
    \vspace{1cm}
\end{titlepage}



\newpage
%------------------------Versioni
\begin{center}
    \textbf{Versioni}
    \\
    \vspace{0.3cm}
    \begin{tabular}{|c|c|c|c|c|}
        \hline
        \textbf{Ver} & \textbf{Data} & \textbf{Redattore} & \textbf{Verificatore} & \textbf{descrizione}\\
        \hline
        0.1 & 05/11/2024 & Gabriele Di Pietro & Filippo Rizzolo & Prima stesura del documento \\
        \hline
    \end{tabular}
\end{center}
%----------------

\printglossary[style=myaltlistgroup,title=]

\end{document}
